\section{Method}
\label{sec:method}

The mesh evaluates Eq.~\eqref{eq:Fc} by a discrete deposit, a sampled kernel,
and a free-space FFT convolution. Let $N$ be the bunch-box nodes per axis,
$M=\mathrm{pad}\cdot N$ the padded FFT grid per axis, $h$ the node spacing in
cell units (node $i$ sits at $i h\,\mathbf{cell}$), and $\mathbf{cell}=(dx,dy,dz)$.
All precision-critical work is in double \citep{SB1,S1-spec}.

\subsection{Deposit}

Each source charge is deposited onto the zero-padded $M^3$ grid with the TSC
shape, at the cell-unit spline argument $W((\mathbf{x}_i-\mathbf{x}_p)/\mathbf{cell})$,
which is correct for any $h$ \citep{S1-spec}. The result
$\rho[i]=\sum_p q_p\,W(i h-\mathbf{p}_{\mathrm{cell}})$ lives in the $N^3$
sub-box of the padded grid, indexed with $M$-stride; the rest is zero. A
fixed reduction order is required for MPI-layout invariance; the single-rank
path uses \texttt{atomicAdd} \citep{S1-spec}.

\subsection{Kernel sampling}

The six TSC-convolved kernels are sampled on the same padded $M^3$ grid, one
thread per FFT point, with the centered (negative-wrapped) index
$d_a = a$ for $a<M/2$ else $a-M$. For each of the six components $c$, the
kernel is evaluated at the physical separation
\begin{equation}
\mathbf{R} = (d_a\,h + \delta_{c,a})\,\mathbf{cell},
\end{equation}
where $\delta_c$ is the Yee staggered offset (Section~\ref{sec:readout}). The
kernel is therefore sampled at the same node spacing $h$ as the deposit, with
the shape one grid cell of its own axis wide
\citep{N2b}. Non-finite samples are set to zero to avoid $0\times\infty$ in
the FFT \citep{S1-spec}.

\subsection{FFT convolution and normalization}

The convolution is a free-space, zero-padded FFT convolution. The vendor plan
is built once, unnormalized, and reused for all thirteen transforms: one
forward for $\rho$, six forwards for the kernels, and six inverse
transformations \citep{S1-spec}. The unnormalized round trip is scaled by
$M^3$; the code divides by $h^3 M^3$ in the read-out
\citep{S1-spec}.

\paragraph{Padding.}
The grid must be zero-padded to $\mathrm{pad}\ge 2$. With $\mathrm{pad}=1$ the
cyclic convolution wraps and a $9.275\times10^{-3}$ periodic-image artefact
appears for a single stationary particle in a $20^3$ box
\citep{N1b}. At $\mathrm{pad}=2$ the result agrees with $\mathrm{pad}=3$ to
$8.4\times10^{-16}$ and with a direct non-cyclic sum to $2.8\times10^{-17}$
\citep{N1b,SB1}. At $\mathrm{pad}\ge2$ the error is at machine precision
($\le 10^{-15}$) \citep{N1b}. The artefact grows for beamed sources: at
$\mathrm{pad}=1$ the $100$-source $\gamma=10$ cloud shows a
$7.5\times10^{-1}$ worst-component error \citep{N1b}.

\paragraph{DC and Nyquist.}
For the R2C/C2R half-spectra the DC bin and the seven self-conjugate even
Nyquist corners must be exactly real; the intermediate Nyquist planes are
Hermitian-paired, not real \citep{N1b}. The R2C/C2R path matches the C2C path
to $1.665\times10^{-16}$ over all six components \citep{N1b}.

\paragraph{The $h^3$ factor.}
The sampled convolution evaluates the continuum integral of
Eq.~\eqref{eq:Fc} by a Riemann sum whose nodes are spaced $h$ cells apart, so
each sample covers a physical volume $h^3$ (in cell units); the read-out
therefore applies the factor
$h_x h_y h_z/(M_x M_y M_z)$ \citep{S1-spec}.

\subsection{Read-out}
\label{sec:readout}

The six field components live at the staggered Yee positions, in cell units
\citep{S1-spec}:
\begin{equation}
\begin{array}{llllll}
E_x & (0.5,0,0) & E_y & (0,0.5,0) & E_z & (0,0,0.5)\\
B_x & (0,0.5,0.5) & B_y & (0.5,0,0.5) & B_z & (0.5,0.5,0).
\end{array}
\end{equation}
The kernel is sampled at these offsets, so the convolution is already
evaluated on the Yee grid. The read-out compacts the $M^3$ FFT result to the
$N^3$ bunch box and applies the $h^3/M^3$ scale and the SI factor
(Section~\ref{sec:fielddef}). No further interpolation is needed
\citep{S1-spec}.

\subsection{Scale invariance}
\label{sec:scaleinvariance}

The kernel is homogeneous of degree $-2$, the deposit shape width equals the
node spacing, and the Yee offsets are in cell units, so every relative metric
depends only on $r/h$, the distance in cell units \citep{N2b}. At a fixed
distance in cells the error is therefore independent of $h$: refining $h$
with the distance held in cells changes nothing, and the measured relative
errors at $h=1,\,1/2,\,1/4,\,1/8$ are identical to machine precision for
$\gamma\in\{1,10,100\}$, aspect 1 and 4 \citep{N2b}. At a fixed physical
distance $r$, refining $h$ increases $r/h$ and does reduce the error, at the
cost of $M^{3}$ FFT work in $1/h$. The only knob is $r/h$: reaching $10^{-3}$
perpendicular at $\gamma=10$ for $r=8$ cells needs $h\approx1/29$,
$2.4\times10^{4}$ times the $M^{3}$ nodes of the $h=1$ run
\citep{N2b,D10}.
